\subsection{Coefficient multiples of horizontal classes}

\begin{lemma}[A-H01: Proved; a uniform coefficient-tail annihilator]
\label{AH01}
Let $E$ be a simplicial bigraded differential algebra over
$S=\Gamma_k(T)$, with coefficient elements constant simplicial
and of form degree zero. Suppose $N(E^{a,0})$ is exact in
homological degree zero for every $a\geq W$, where $W\geq1$.
Then the ideal
$I_W=\bigoplus_{a\geq W}kT^{[a]}$ annihilates
\[
 H_n\bigl(N(E^{*,q})\bigr)
\]
for every $q,n\geq0$. This conclusion uses no vertical
exactness, flatness, or divided-power operation on the
positive-form part of $E$.
\end{lemma}

\begin{proof}
Use normalization
$N_n(V)=\bigcap_{j=1}^{n}\ker(d_j:V_n\to V_{n-1})$
and differential $d_0$. Fix $a\geq W$ and a normalized
cycle $c\in N_n(E^{w,q})$. Exactness in homological
degree zero gives
$h\in N_1(E^{a,0})$ with
$d_0h=T^{[a]}$ and $d_1h=0$.
For $0\leq i\leq n$, let
\[
 H_i=(s_n\cdots\widehat{s_i}\cdots s_0 h)(s_i c)
       \in E_{n+1}^{a+w,q},
 \qquad H=\sum_{i=0}^{n}(-1)^iH_i .
\]
The degeneracies on the first factor are written in
decreasing order; exactly $n$ of them occur. For $n=0$
this means $H=h(s_0c)$.

We check the normalized identities directly. If
$1\leq j\leq n$, the simplicial identities and the
equations $d_jc=0$ show that $d_jH_i=0$ except possibly
when $i=j-1$ or $i=j$. For these two terms,
$d_j s_{j-1}c=d_j s_jc=c$, and the degeneracies on
$h$ reduce to the same expression. Their signs are
opposite, so $d_jH=0$.
For $j=n+1$ only $i=n$ can contribute; its first
factor is a degeneracy of $d_1h=0$.
Thus $H\in N_{n+1}(E^{a+w,q})$.

For $i>0$, the equality
$d_0s_i c=s_{i-1}d_0c=0$ kills $d_0H_i$.
For $i=0$, the first factor of $d_0H_0$ is the
constant simplex $T^{[a]}$, because it is the
iterated degeneracy of $d_0h$; its second factor is
$d_0s_0c=c$. Therefore
\[
 d_0H=T^{[a]}c.
\]
This proves that each coefficient basis element in
$I_W$ acts by zero on each horizontal homology group.
Linearity gives the ideal-annihilator assertion.
\end{proof}

In particular, the equations asserting that coefficient-tail
multiples of a closed horizontal cycle have primitives
are already satisfied by the displayed formulas. These
equations alone do not establish that the original class
is zero. A use of the scaled bar map of A-L05 in a
horizontal repair must still give a proof that the
resulting attachment fills the actual unscaled class;
the lemma makes no assertion about other relations or
operations introduced by such an attachment.
